\subsection{Überwachung einer Spannung mit LEDs}
% Alt: Aufgabe 1
\Aufgabe{
    \label{Aufg1}
    Gegeben ist die abgebildete Schaltung, mit dieser soll die Spannung $U_\mathrm{a}$ überwacht werden. $U_\mathrm{a}$ darf dabei nur im folgenden Bereich liegen $\mathrm{4,75\,V} < U_\mathrm{a} < \mathrm{5,25\,V}$. Wenn die Spannung über oder unterschritten ist, soll jeweils
    eine der LEDs aufleuchten. Erzgänzen Sie die Schaltung.
    \begin{figure}[H]
        \centering
        \input{Tikz/src/FigSpannungsueberwachungLED.tex}\alttext{Das Bild zeigt ein elektrisches Schaltbild mit zwei Operationsverstärkern, die mit N1 (oben) und N2 (unten) bezeichnet sind. Beide Operationsverstärker werden mit einer symmetrischen Versorgung von +15 V und -15 V betrieben, die links im Schaltbild eingezeichnet ist und jeweils mit den Anschlüssen U+ und U- verbunden sind.
            Von links unten ist eine Spannung \(U_\mathrm{a}\) mit einem Pfeil nach unten dargestellt. Diese Spannung wird zu beiden Operationsverstärkern geführt und liegt an deren Eingängen an. Jeder Operationsverstärker besitzt rechts einen Ausgang. Am Ausgang von N1 ist eine LED mit der Bezeichnung V1 angeschlossen, am Ausgang von N2 eine LED mit der Bezeichnung V2. Die jeweils andere Seite der LEDs ist mit -15 V verbunden.
            Unten links ist zusätzlich ein Masseanschluss eingezeichnet. Das Schaltbild stellt somit zwei parallel betriebene Operationsverstärker mit gemeinsamer Eingangsspannung und jeweils einer LED am Ausgang dar, die das jeweilige Ausgangssignal anzeigen.}
        \label{fig:FigSpannungsueberwachungLED}
    \end{figure}

}

\Loesung{

    \begin{figure}[H]
        \centering
        \input{Tikz/src/LsgSpannungsueberwachungLED1.tex}\alttext{Die Abbildung zeigt eine Schaltung mit zwei Operationsverstärkern, die übereinander angeordnet sind und jeweils mit einer symmetrischen Versorgung von \(+15 \ V\) und \(-15 \ V\) betrieben werden.
            Oben ist der Operationsverstärker N1 dargestellt, darunter der Operationsverstärker N2. Beide besitzen einen invertierenden Eingang (−), einen nichtinvertierenden Eingang (+) sowie Anschlüsse für die positive und negative Versorgungsspannung \(U+\) und \(U-\).Die Versorgungsspannungen werden links eingespeist und zu beiden Verstärkern weitergeführt.
            Die Ausgänge beider Operationsverstärker befinden sich jeweils auf der rechten Seite. An jedem Ausgang ist eine LED angeschlossen (V1 bei N1 und V2 bei N2), die jeweils gegen \(-15 \ V\). geschaltet ist. Dadurch leuchtet die entsprechende LED abhängig vom Ausgangspotential des jeweiligen Operationsverstärkers.
            Links unten ist die Ausgangsspannung \(U_\mathrm{a}\) eingezeichnet, die gegen Masse gemessen wird.
            Im unteren linken Bereich ist zusätzlich ein Spannungsteiler aus den Widerständen \(R_1\) und \(R_2\) eingezeichnet. Dieser liegt zwischen einer positiven Versorgungsspannung und Masse. Über dem Widerstand \(R_1\) fällt eine Spannung von \(5,25 \ \mathrm{V}\) ab, die in der Abbildung markiert ist. Dieser Spannungspunkt ist mit dem oberen Operationsverstärker verbunden.}
        \label{fig:LsgSpannungsueberwachungLED1}
    \end{figure}
    \begin{equation*}
        \begin{aligned}
            \frac{U_1}{U_0}                   & = \frac{R_1}{R_1 + R_2}        \\
            \frac{U_1 \cdot (R_1 + R_2)}{U_0} & = R_1                          \\
            U_1\cdot (R_1+R_2)                & = R_1 \cdot U_0                \\
            U_1\cdot R_1 + U_1\cdot R_2       & = R_1 \cdot U_0                \\
            U_1\cdot R_1                      & = R_1 \cdot U_0 - U_1\cdot R_2 \\
            U_1\cdot R_2                      & = R_1 \cdot (U_0 - U_1)        \\
            \frac{U_1 \cdot R_2}{U_0 - U_1}   & = R_1
        \end{aligned}
    \end{equation*}

    \[
        R_1 = 268.96 \, \Omega
    \]

    \begin{figure}[H]
        \centering
        \input{Tikz/src/LsgSpannungsueberwachungLED2.tex}\alttext{Die Abbildung zeigt eine Schaltung mit zwei Operationsverstärkern, die übereinander angeordnet sind und jeweils mit einer symmetrischen Versorgung von \(+15 \ \mathrm{V}\) und \(-15 \ \mathrm{V}\) betrieben werden.
            Oben ist der Operationsverstärker N1 dargestellt, darunter der Operationsverstärker N2. Beide besitzen einen invertierenden Eingang (-), einen nichtinvertierenden Eingang (+) sowie Anschlüsse für die positive und negative Versorgungsspannung \(U+\) und \(U-\).Die Versorgungsspannungen werden links eingespeist und zu beiden Verstärkern weitergeführt.
            Die Ausgänge beider Operationsverstärker befinden sich jeweils auf der rechten Seite. Am Ausgang von N1 ist über den Widerstand \(R_1\) eine LED (V1) angeschlossen, die gegen \(-15 \ \mathrm{V}\) geschaltet ist. Am Ausgang von N2 ist über den Widerstand \(R_2\) ebenfalls eine LED (V2) angeschlossen, die ebenfalls gegen \(-15 \ \Mathrm{V}\) geschaltet ist. Dadurch zeigt jede LED abhängig vom Ausgangspotential des jeweiligen Operationsverstärkers den entsprechenden Schaltzustand an.
            Links unten ist die Spannung \(U_\mathrm{a}\) eingezeichnet, die gegen Masse gemessen wird. Diese Spannung wird auf beide Operationsverstärker geführt.
            Im linken unteren Bereich sind zwei ideale Spannungsquellen eingezeichnet, die als Referenzspannungen dienen. Eine Referenzspannung von \(5,25 \ \mathrm{V}\) ist mit dem oberen Operationsverstärker N1 verbunden. Eine zweite Referenzspannung von \(4,75 \ \mathrm{V}\) ist mit dem unteren Operationsverstärker N2 verbunden. Beide Referenzspannungen sind jeweils gegen Masse bezogen.}
        \label{fig:LsgSpannungsueberwachungLED2}
    \end{figure}

    \begin{figure}[H]
        \centering
        \input{Tikz/src/LsgSpannungsueberwachungLED3.tex}\alttext{Die Abbildung zeigt eine Schaltung mit zwei Operationsverstärkern, die übereinander angeordnet sind und jeweils mit einer symmetrischen Versorgung von \(+15 \ \mathrm{V}\) und \(-15 \ \mathrm{V}\) betrieben werden.
            Oben ist der Operationsverstärker N1 dargestellt, darunter der Operationsverstärker N2. Beide besitzen einen invertierenden Eingang (-), einen nichtinvertierenden Eingang (+) sowie Anschlüsse für die positive und negative Versorgungsspannung \(U+\) und \(U-\).Die Versorgungsspannungen werden links eingespeist und zu beiden Verstärkern weitergeführt.
            Die Ausgänge beider Operationsverstärker befinden sich jeweils auf der rechten Seite. Am Ausgang von N1 ist über den Widerstand \(R_4\) eine LED (V1) angeschlossen, die gegen \(-15 \ V\) geschaltet ist. Für diese LED ist ein Strom von etwa \(2 \ \mathrm{mA}\) sowie eine Durchlassspannung von ungefähr \(2 \ \mathrm{V}\)angegeben. Am Ausgang von N2 ist über den Widerstand \(R_5\) ebenfalls eine LED (V2) angeschlossen, die ebenfalls gegen \(-15 \ \mathrm{V}\) geschaltet ist. Dadurch zeigt jede LED abhängig vom jeweiligen Ausgangspotential des Operationsverstärkers den Schaltzustand an.
            Links unten ist die Spannung \(U_\mathrm{a}\) eingezeichnet, die gegen Masse gemessen wird. Diese Spannung wird auf beide Operationsverstärker geführt.
            Im linken Bereich der Schaltung befindet sich ein dreistufiger Spannungsteiler aus den Widerständen \(R_1\), \(R_2\) und \(R_3\).Dieser liegt zwischen \(+15 \ V\) und Masse. Über \(R_1\) ällt eine Spannung von \(5,75 \ \mathrm{V}\)ab, über \(5,75 \ \mathrm{V}\) ab, über \(R_2\) ine Spannung von \(o,5 \ \mathrm{V}\) und über \(R_3\) eine Spannung von \(4,75 \ \mathrm{V}\).
            Der Zwischenpunkt zwischen \(R_1\) und \(R_2\) ist mit dem oberen Operationsverstärker N1 verbunden und stellt die obere Referenzspannung dar. Der Zwischenpunkt zwischen \(R_2\) und \(R_3\) ist mit dem unteren Operationsverstärker N2 verbunden und bildet die untere Referenzspannung.}
        \label{fig:LsgSpannungsueberwachungLED3}
    \end{figure}
    \begin{align*}
        R_1 & = 3.75 \,\text{k}\Omega                                                   \\
        R_2 & = 500 \,\Omega                                                            \\
        R_3 & = 4.75 \,\text{k}\Omega                                                   \\
        R_4 & = \frac{28 \, \mathrm{V}}{20 \, \mathrm{mA}} \approx 1,4 \,\text{k}\Omega
    \end{align*}
    Siehe: Abschnitt \ref{fig: Verschaltung Verstaerker} und Tabelle \ref{tab:Grundschaltungen}
}